% Figure: cross-section of a microstrip transmission line. A narrow top
% conductor of width w sits on a dielectric substrate of thickness h over a
% full ground plane. Field lines run from the strip down to the ground both
% through the substrate and through the air above (fringing), which is why
% the effective permittivity lies between 1 and epsilon_r. Pure tikz drawing.
\begin{figure}[htbp]
\centering
\begin{tikzpicture}[scale=1.2, >=latex]
  % substrate block
  \fill[engblue!8] (-3,0) rectangle (3,1.1);
  \draw[thick] (-3,0) rectangle (3,1.1);
  \node[font=\scriptsize, anchor=west] at (3.05,0.55)
    {substrate $\varepsilon_r$};
  % ground plane (thick bar at bottom)
  \fill[enggray] (-3,-0.18) rectangle (3,0);
  \node[font=\scriptsize, anchor=west] at (3.05,-0.09) {ground plane};
  % top strip conductor of width w, centred
  \fill[engred] (-0.8,1.1) rectangle (0.8,1.28);
  \node[engred, font=\scriptsize, anchor=south] at (0,1.30) {strip, width $w$};
  % substrate thickness marker h
  \draw[<->, enggray] (-2.7,0) -- (-2.7,1.1);
  \node[enggray, font=\scriptsize, anchor=east] at (-2.72,0.55) {$h$};
  % field lines from strip edges down to the ground: substrate paths
  \draw[engblue, ->] (-0.4,1.1) .. controls (-0.4,0.6) .. (-0.4,0);
  \draw[engblue, ->] (0.4,1.1) .. controls (0.4,0.6) .. (0.4,0);
  \draw[engblue, ->] (0,1.1) -- (0,0);
  % fringing field lines curling out through the air above the substrate
  \draw[engamber, ->] (-0.8,1.19) .. controls (-1.7,1.0) and (-1.7,0.2)
    .. (-1.2,0);
  \draw[engamber, ->] (0.8,1.19) .. controls (1.7,1.0) and (1.7,0.2)
    .. (1.2,0);
  \node[engamber, font=\scriptsize, anchor=south west] at (1.2,0.5)
    {fringing (air)};
\end{tikzpicture}
\caption[Cross-section of a microstrip line and its fringing field]{A
microstrip line: a narrow top conductor of width $w$ printed on a
dielectric substrate of thickness $h$ over a solid ground plane. Most of
the field runs straight down through the substrate, but a substantial part
fringes out through the air above the strip, so the wave samples a mixture
of the two media. The effective permittivity that sets the guided
wavelength lies between $1$ and $\varepsilon_r$, closer to $\varepsilon_r$
for a wide strip on a thin substrate where the field is well confined, and
lower for a narrow strip where more field escapes into the air. The same
fringing at the open ends of a printed patch is what makes the patch
resonate slightly longer than its metal, the fringing-length correction of
the microstrip-patch case study.}
\label{fig:vol04ch10:microstrip-cross-section}
\end{figure}
